\chapter{Excell Una fantastica Tabla de Componentes}
Ahora que sabes como crear custom React Componentes y muchas cosas mas. Pongamos todo esto a prueba creando una tabla de datos\\
\section{Primero los Datos}
Aqui tenemos un ejemplo
\begin{lstlisting}[language=html]
const headers = ['Book', 'Author', 'Language', 'Published', 'Sales'];
const data = [
[
'A Tale of Two Cities', 'Charles Dickens',
'English', '1859', '200 million',
],
[
'Le Petit Prince (The Little Prince)', 'Antoine de Saint-Exupéry',
'French', '1943', '150 million',
],
[
"Harry Potter and the Philosopher's Stone", 'J. K. Rowling',
'English', '1997', '120 million',
],
[
'And Then There Were None', 'Agatha Christie',
'English', '1939', '100 million',
],
[
'Dream of the Red Chamber', 'Cao Xueqin',
'Chinese', '1791', '100 million',
],
[
'The Hobbit', 'J. R. R. Tolkien',
'English', '1937', '100 million',
],
];
\end{lstlisting}
\section{Table Headers Loop}
Creamos el iterador para contruir la cabecera
\begin{lstlisting}[language=html]
<!doctype html>
<html>
  <header>
    <meta charset="utf-8" />
    <title>Header Loop Table</title>
    <link rel="stylesheet" type="text/css" href="03.table.css" />
  </header>
  <body>
    <div id="app"></div>
    <script src="react/react.js"></script>
    <script src="react/react-dom.js"></script>
    <script src="react/babel.js"></script>
    <script type="text/babel">
      const headers = ["Book", "Author", "Language", "Published", "Sales"];
      const data = [
        [
          "A Tale of Two Cities",
          "Charles Dickens",
          "English",
          "1859",
          "200 million",
        ],
        [
          "Le Petit Prince (The Little Prince)",
          "Antoine de Saint-Exupéry",
          "French",
          "1943",
          "150 million",
        ],
        [
          "Harry Potter and the Philosopher's Stone",
          "J. K. Rowling",
          "English",
          "1997",
          "120 million",
        ],
        [
          "And Then There Were None",
          "Agatha Christie",
          "English",
          "1939",
          "100 million",
        ],
        [
          "Dream of the Red Chamber",
          "Cao Xueqin",
          "Chinese",
          "1791",
          "100 million",
        ],
        ["The Hobbit", "J. R. R. Tolkien", "English", "1937", "100 million"],
      ];
      class Excell extends React.Component {
        render() {
          const headers = [];
          for (const title of this.props.headers) {
            headers.push(<th>{title}</th>);
          }
          return (
            <table>
              <thead>
                <tr>{headers}</tr>
              </thead>
            </table>
          );
        }
      }
      ReactDOM.render(
        <Excell headers={headers} />,
        document.getElementById("app"),
      );
    </script>
  </body>
</html>
\end{lstlisting}
Como puedes ver renderisamos el array headers en tr el cual fue compuesto por el metodo for 
\section{Agregando Content $<td>$}
Ahora que tiene un cabecera de tabla bien es hora de agregar el body, ahora renderisaremos un array de dos dimenciones filas y columnas \\
Para pasar el dato a Excel usaremo un nuevo prop llamado initialData. Porque initial y no lo sata? Es por manejo esperado. El usuario podra editarlo ordenarlo y un poco mas, en otras palabras el estado del componente cambiara Entonces usemos this.state.data para mantener el trackeo de los cambios y usar this.props.initialData
\begin{lstlisting}[language=html]
<!doctype html>
<html>
  <head>
    <title>Sort Table</title>
    <meta charset="utf-8" />
    <link rel="stylesheet" type="text/css" href="03.table.css" />
  </head>
  <body>
    <div id="app"></div>
    <script src="react/react.js"></script>
    <script src="react/react-dom.js"></script>
    <script src="react/babel.js"></script>
    <script type="text/babel">
      const headers = ["Book", "Author", "Language", "Published", "Sales"];
      const data = [
        [
          "A Tale of Two Cities",
          "Charles Dickens",
          "English",
          "1859",
          "200 million",
        ],
        [
          "Le Petit Prince (The Little Prince)",
          "Antoine de Saint-Exupéry",
          "French",
          "1943",
          "150 million",
        ],
        [
          "Harry Potter and the Philosopher's Stone",
          "J. K. Rowling",
          "English",
          "1997",
          "120 million",
        ],
        [
          "And Then There Were None",
          "Agatha Christie",
          "English",
          "1939",
          "100 million",
        ],
        [
          "Dream of the Red Chamber",
          "Cao Xueqin",
          "Chinese",
          "1791",
          "100 million",
        ],
        ["The Hobbit", "J. R. R. Tolkien", "English", "1937", "100 million"],
      ];
....
      class Excell extends React.Component {
        constructor(props) {
          super();
          this.state = { data: props.initialData };
        }
        render() {
          const headers = [];
          for (const idx in this.props.headers) {
            const header = this.props.headers[idx];
            headers.push(<th>{header}</th>);
          }
          return (
            <div>
              <h3>Table</h3>
              <table>
                <thead>
                  <tr>{headers}</tr>
                </thead>
                <tbody>
                  {this.state.data.map((row, idx) => (
                    <tr key={idx}>
                      {row.map((cell, idx) => (
                        <td key={idx}>{cell}</td>
                      ))}
                    </tr>
                  ))}
                </tbody>
              </table>
            </div>
          );
        }
      }
      ReactDOM.render(
        <Excell headers={headers} initialData={data} />,
        document.getElementById("app"),
      );
    </script>
  </body>
</html>
\end{lstlisting}
En este caso se usara un doble map para hacer el recorrido de cada celda en cada fila y columna tambien podria usarse un doble for anidado pero no seria buena practica
\section{Prop Types}
La habilidad para especificar tipos de variables como (string, numeros, booleans etc) no existe en Javacript. Pero los desarrolladores de otros lenguajes, y estos que trabajan en proyecto grandescon muchos otros desarrolladores, lo olvidan.\\
Dos opciones populares existe nque permiten escribir JavaScript con tipos: Flow y TypeScript. Puedes esar estos para escribir applicaciones React. \\
Pero otra opcion existente cual es limitada para solo especificar el tipo de props en tu componente con prop types. Ellos fueron parte de React, inicialmente pero fueron movidos a una librearia separada en React v15.5\\
Prop types te permite ser mas especifico como que dato de la tabla tomar y como un resultado un error en el developer console aparecera. Puedes configurar el prop type como esto \\
Para esto necesitaras una nueva libreria 
\begin{lstlisting}[language=bash]
curl -L https://unpkg.com/prop-types/prop-types.js > ~/reactbook/react/prop-types.js
\end{lstlisting}
y especificar lo siguiente
\begin{lstlisting}[language=html]
<!doctype html>
<html>
  <head>
    <title>Sort Table</title>
    <meta charset="utf-8" />
    <link rel="stylesheet" type="text/css" href="03.table.css" />
  </head>
  <body>
    <div id="app"></div>
    <script src="react/react.js"></script>
    <script src="react/react-dom.js"></script>
    <script src="react/babel.js"></script>
    <script src="react/prop-types.js"></script>
    <script type="text/babel">
      const headers = ["Book", "Author", "Language", "Published", "Sales"];
      const data = [
        [
          "A Tale of Two Cities",
          "Charles Dickens",
          "English",
          "1859",
          "200 million",
        ],
        [
          "Le Petit Prince (The Little Prince)",
          "Antoine de Saint-Exupéry",
          "French",
          "1943",
          "150 million",
        ],
        [
          "Harry Potter and the Philosopher's Stone",
          "J. K. Rowling",
          "English",
          "1997",
          "120 million",
        ],
        [
          "And Then There Were None",
          "Agatha Christie",
          "English",
          "1939",
          "100 million",
        ],
        [
          "Dream of the Red Chamber",
          "Cao Xueqin",
          "Chinese",
          "1791",
          "100 million",
        ],
        ["The Hobbit", "J. R. R. Tolkien", "English", "1937", "100 million"],
      ];
....
      class Excell extends React.Component {
        constructor(props) {
          super();
          this.state = { data: props.initialData };
        }
        render() {
          const headers = [];
          for (const idx in this.props.headers) {
            const header = this.props.headers[idx];
            headers.push(<th key={idx}>{header}</th>);
          }
          return (
            <div>
              <h3>Table</h3>
              <table>
                <thead>
                  <tr>{headers}</tr>
                </thead>
                <tbody>
                  {this.state.data.map((row, idx) => (
                    <tr key={idx}>
                      {row.map((cell, idx) => (
                        <td key={idx}>{cell}</td>
                      ))}
                    </tr>
                  ))}
                </tbody>
              </table>
            </div>
          );
        }
      }
      Excell.propTypes = {
        header: PropTypes.arrayOf(PropTypes.string),
        initialData: PropTypes.arrayOf(PropTypes.arrayOf(PropTypes.string)),
      };
      ReactDOM.render(
        <Excell headers={headers} initialData={data} />,
        document.getElementById("app"),
      );
    </script>
  </body>
</html>
\end{lstlisting}
\section{Sorting}
Cuantas veces viste una tabla en la  que puedes ordenarla diferente? Es trivial hacer esto con React. Acualmente este es un ejemplo donde React porque todo lo que necesitas hacer es ordenar el array y todo la UI lo actualizara por ti\\
Por conveniencia y leible todos los metodos de ordenamiento estan en un metodo sort()
\begin{lstlisting}[language=html]
<!doctype html>
<html>
  <head>
    <title>Sort Table</title>
    <meta charset="utf-8" />
    <link rel="stylesheet" type="text/css" href="03.table.css" />
  </head>
  <body>
    <div id="app"></div>
    <script src="react/react.js"></script>
    <script src="react/react-dom.js"></script>
    <script src="react/babel.js"></script>
    <script src="react/prop-types.js"></script>
    <script type="text/babel">
      const headers = ["Book", "Author", "Language", "Published", "Sales"];
      const data = [
        [
          "A Tale of Two Cities",
          "Charles Dickens",
          "English",
          "1859",
          "200 million",
        ],
        [
          "Le Petit Prince (The Little Prince)",
          "Antoine de Saint-Exupéry",
          "French",
          "1943",
          "150 million",
        ],
        [
          "Harry Potter and the Philosopher's Stone",
          "J. K. Rowling",
          "English",
          "1997",
          "120 million",
        ],
        [
          "And Then There Were None",
          "Agatha Christie",
          "English",
          "1939",
          "100 million",
        ],
        [
          "Dream of the Red Chamber",
          "Cao Xueqin",
          "Chinese",
          "1791",
          "100 million",
        ],
        ["The Hobbit", "J. R. R. Tolkien", "English", "1937", "100 million"],
      ];
      function clone(o) {
        return JSON.parse(JSON.stringify(o));
      }
      class Excell extends React.Component {
        constructor(props) {
          super();
          this.state = { data: props.initialData };
          this.sort = this.sort.bind(this);
        }
        sort(e) {
          const column = e.target.cellIndex;
          console.log(column, e.target);
          const data = clone(this.state.data);
          data.sort((a, b) => {
            if (a[column] === b[column]) {
              return 0;
            }
            return a[column] > b[column] ? 1 : -1;
          });
          this.setState({ data });
        }
        render() {
          const headers = [];
          for (const idx in this.props.headers) {
            const header = this.props.headers[idx];
            headers.push(<th key={idx}>{header}</th>);
          }
          return (
            <div>
              <h3>Table</h3>
              <table>
                <thead onClick={this.sort}>
                  <tr>{headers}</tr>
                </thead>
                <tbody>
                  {this.state.data.map((row, idx) => (
                    <tr key={idx}>
                      {row.map((cell, idx) => (
                        <td key={idx}>{cell}</td>
                      ))}
                    </tr>
                  ))}
                </tbody>
              </table>
            </div>
          );
        }
      }
      Excell.propTypes = {
        headers: PropTypes.arrayOf(PropTypes.string),
        initialData: PropTypes.arrayOf(PropTypes.arrayOf(PropTypes.string)),
      };
      ReactDOM.render(
        <Excell headers={headers} initialData={data} />,
        document.getElementById("app"),
      );
    </script>
  </body>
</html>
\end{lstlisting}
Aqui usamos el metodo sort el cual puede desmenusar el array en y empezara comprar una fila con la siguiente y devolver 1, -1, 0 si se necesita ordenarlo
\section{Sorting UI Cues}
La tabla es ordenada sin problemas pero no es claro que columna fue ordenada, actualicemos la UI para mostrar una flecha basada en la column
Para mantener el tracking de los nuevos estados necesitaremos nuevos estados\\
this.state.sortby, Nombre de la columna que esta actualmente ordenando la tabla
this.state.descending, True or False si esta ordenada ascendiente o descendiente
\begin{lstlisting}[language=html]
<!doctype html>
<html>
  <head>
    <title>Sort Table</title>
    <meta charset="utf-8" />
    <link rel="stylesheet" type="text/css" href="03.table.css" />
  </head>
  <body>
    <div id="app"></div>
    <script src="react/react.js"></script>
    <script src="react/react-dom.js"></script>
    <script src="react/babel.js"></script>
    <script src="react/prop-types.js"></script>
    <script type="text/babel">
      const headers = ["Book", "Author", "Language", "Published", "Sales"];
      const data = [
        [
          "A Tale of Two Cities",
          "Charles Dickens",
          "English",
          "1859",
          "200 million",
        ],
        [
          "Le Petit Prince (The Little Prince)",
          "Antoine de Saint-Exupéry",
          "French",
          "1943",
          "150 million",
        ],
        [
          "Harry Potter and the Philosopher's Stone",
          "J. K. Rowling",
          "English",
          "1997",
          "120 million",
        ],
        [
          "And Then There Were None",
          "Agatha Christie",
          "English",
          "1939",
          "100 million",
        ],
        [
          "Dream of the Red Chamber",
          "Cao Xueqin",
          "Chinese",
          "1791",
          "100 million",
        ],
        ["The Hobbit", "J. R. R. Tolkien", "English", "1937", "100 million"],
      ];
      function clone(o) {
        return JSON.parse(JSON.stringify(o));
      }
      class Excell extends React.Component {
        constructor(props) {
          super();
          this.state = {
            data: props.initialData,
            sortby: null,
            descending: false,
          };
          this.sort = this.sort.bind(this);
        }
        sort(e) {
          const column = e.target.cellIndex;
          const data = clone(this.state.data);
          const descending =
            this.state.sortby == column && !this.state.descending;
          data.sort((a, b) => {
            if (a[column] === b[column]) {
              return 0;
            }
            return descending
              ? a[column] < b[column]
                ? 1
                : -1
              : a[column] > b[column]
                ? 1
                : -1;
          });
          this.setState({
            data,
            sortby: column,
            descending,
          });
        }
        render() {
          const headers = [];
          for (const idx in this.props.headers) {
            const header = this.props.headers[idx];
            headers.push(<th key={idx}>{header}</th>);
          }
          return (
            <div>
              <h3>Table</h3>
              <table>
                <thead onClick={this.sort}>
                  <tr>
                    {this.props.headers.map((title, idx) => {
                      if (this.state.sortby === idx) {
                        title += this.state.descending ? " \u2191" : " \u2193";
                      }
                      return <th key={idx}>{title}</th>;
                    })}
                  </tr>
                </thead>
                <tbody>
                  {this.state.data.map((row, idx) => (
                    <tr key={idx}>
                      {row.map((cell, idx) => (
                        <td key={idx}>{cell}</td>
                      ))}
                    </tr>
                  ))}
                </tbody>
              </table>
            </div>
          );
        }
      }
      Excell.propTypes = {
        headers: PropTypes.arrayOf(PropTypes.string),
        initialData: PropTypes.arrayOf(PropTypes.arrayOf(PropTypes.string)),
      };
      ReactDOM.render(
        <Excell headers={headers} initialData={data} />,
        document.getElementById("app"),
      );
    </script>
  </body>
</html>
\end{lstlisting}
\section{Editable cell}
Lo primero es configurar un event handler simple. En un doble click el componente sera renombrado
\begin{lstlisting}[language=html]
<!doctype html>
<html>
  <head>
    <title>Sort Table</title>
    <meta charset="utf-8" />
    <link rel="stylesheet" type="text/css" href="03.table.css" />
  </head>
  <body>
    <div id="app"></div>
    <script src="react/react.js"></script>
    <script src="react/react-dom.js"></script>
    <script src="react/babel.js"></script>
    <script src="react/prop-types.js"></script>
    <script type="text/babel">
      const headers = ["Book", "Author", "Language", "Published", "Sales"];
      const data = [
        [
          "A Tale of Two Cities",
          "Charles Dickens",
          "English",
          "1859",
          "200 million",
        ],
        [
          "Le Petit Prince (The Little Prince)",
          "Antoine de Saint-Exupéry",
          "French",
          "1943",
          "150 million",
        ],
        [
          "Harry Potter and the Philosopher's Stone",
          "J. K. Rowling",
          "English",
          "1997",
          "120 million",
        ],
        [
          "And Then There Were None",
          "Agatha Christie",
          "English",
          "1939",
          "100 million",
        ],
        [
          "Dream of the Red Chamber",
          "Cao Xueqin",
          "Chinese",
          "1791",
          "100 million",
        ],
        ["The Hobbit", "J. R. R. Tolkien", "English", "1937", "100 million"],
      ];
      function clone(o) {
        return JSON.parse(JSON.stringify(o));
      }
      class Excell extends React.Component {
        constructor(props) {
          super();
          this.state = {
            data: props.initialData,
            sortby: null,
            descending: false,
            edit: null,
          };
          this.sort = this.sort.bind(this);
          this.showEditor = this.showEditor.bind(this);
          this.save = this.save.bind(this);
        }
        sort(e) {
          const column = e.target.cellIndex;
          const data = clone(this.state.data);
          const descending =
            this.state.sortby == column && !this.state.descending;
          data.sort((a, b) => {
            if (a[column] === b[column]) {
              return 0;
            }
            return descending
              ? a[column] < b[column]
                ? 1
                : -1
              : a[column] > b[column]
                ? 1
                : -1;
          });
          this.setState({
            data,
            sortby: column,
            descending,
          });
        }
        save(e) {
          e.preventDefault();
          const input = e.target.firstChild;
          const data = clone(this.state.data);
          data[this.state.edit.row][this.state.edit.column] = input.value;
          this.setState({ edit: null, data });
        }
        showEditor(e) {
          this.setState({
            edit: {
              row: parseInt(e.target.parentNode.dataset.row, 10),
              column: e.target.cellIndex,
            },
          });
        }
        render() {
          const headers = [];
          for (const idx in this.props.headers) {
            const header = this.props.headers[idx];
            headers.push(<th key={idx}>{header}</th>);
          }
          return (
            <div>
              <h3>Table</h3>
              <table>
                <thead onClick={this.sort}>
                  <tr>
                    {this.props.headers.map((title, idx) => {
                      if (this.state.sortby === idx) {
                        title += this.state.descending ? " \u2191" : " \u2193";
                      }
                      return <th key={idx}>{title}</th>;
                    })}
                  </tr>
                </thead>
                <tbody onDoubleClick={this.showEditor}>
                  {this.state.data.map((row, rowidx) => (
                    <tr key={rowidx} data-row={rowidx}>
                      {row.map((cell, columnidx) => {
                        const edit = this.state.edit;
                        if (
                          edit &&
                          edit.row == rowidx &&
                          edit.column == columnidx
                        ) {
                          cell = (
                            <form onSubmit={this.save}>
                              <input type="text" defaultValue={cell} />
                            </form>
                          );
                        }
                        return <td key={columnidx}>{cell}</td>;
                      })}
                    </tr>
                  ))}
                </tbody>
              </table>
            </div>
          );
        }
      }
      Excell.propTypes = {
        headers: PropTypes.arrayOf(PropTypes.string),
        initialData: PropTypes.arrayOf(PropTypes.arrayOf(PropTypes.string)),
      };
      ReactDOM.render(
        <Excell headers={headers} initialData={data} />,
        document.getElementById("app"),
      );
    </script>
  </body>
</html>
\end{lstlisting}
\section{Conclusiones y Virtual DOM Diffs}
En este punto el feature edit esta completo. No tomo mucho codigo
\begin{itemize}
  \item Mantenimos el tracking de que celda editar via this.state.edit
  \item Renderisamos el input field cuando mostraba la tabla si el row y la columna conincidian 
  \item Actualizamos el array data con el nuevo valor desde el input field
\end{itemize}
Tan pronto como tu setState() con el nuevo dato React llama el componente metodo render() y la UI magicamnete se actualiza. Esto podria verse como no es eficiente hacer render de toda la tabla para solo una celda. Es un hecho React solo actualizara una celda \\
Si abres el dev tool puedes ver que partes del DOM tree son actualizados como tu interaccion con tu aplicacion. puedes ver el dev tools highlighting el DOM cambia despues cambiar\\
Detras de escenas React llama tu metodo render() y crea un lighweight tree representation del DOM result. Esto es conocido como un DOM tree virtual. Cuando el metodo render() es llamado de nuevo (despues una llamada a setSetat())React toma el arbol virtual antes y despues y computa una comparacion Basado en esta comparacion React descrubre las operaciones minimas requeridas DOM(appendChild(), textContent, etc) para llevar los cambios dentro el browser DOM\\
Hay solo un cambio requerido para la celda y no es necesario renderisar toda la tabla. Calculando el minimo set de cambios y agrupando y computando el DOM, React "toca" el DOM ligeramente, como problema conocido que las operaciones DOM son lentas(comparadas a operaciones puras Javacript) llamadas de funciones,etc.\\
Y amenudo hay cuellos de botellas en el rendimiendo de renderisacion de una web application\\
React tiene su espalda cuando viene a llevar a cabo y actualizar la UI
\begin{itemize}
  \item Tocando el DOM ligeramente
  \item Usando delegacion eventos para interaccion de usuarios
\end{itemize}
\section{Search}
A continuacion agregamos un search feature a la tabla 
\begin{enumerate}
  \item Agregar un boton para prender el nuevo feature y luego apagarlo
  \item Si el search esta prendido Agrega un row de input donde cada busqueda en la columna correspondiente
  \item Como un user type en un input box, filtrar el array de state.data para mostrar solo el contenido que coincida
\end{enumerate}
\section{State y UI}
Lo primero a hacer es actualizar el contructor por
\begin{itemize}
  \item Agregando una propiedad search para el objeto this.state para mantener el trackeo de donde se esta haciendo la busqueda
  \item Agreando dos nuevos metodos: this.togglesearch() para activar la cajas y apagarlo y this.search() hara la busqueda
  \item Poniendo una nueva clase  this.preSearchData
  \item Actualizando el dato entrada incicial con un consecutive ID para ayudar a identificar las filas cuando editamos contenidos de datos filtrados
\end{itemize}
La propiedad search permitira saber si el usuario le dio click a mostrar u ocultar la barra de busqueda mediante el metodo this.togglesearch()\\
El metodo search hara la busqueda escuchando el texto escrito \\
Cuando hacemos click en el boton search la funcion toggleSearch() es invocada esta funcion se prende y apaga. Esto es hecho por las tareas
\begin{itemize}
  \item Seteando el this.state.search en true o false
  \item Cuando es true, recuerda el dato actual
  \item Cuando es false revierte al dato 
\end{itemize}
\section{Filtrando el contenido}
El feature search es simple llama al Array.prototype.filter() en el array y retorna un array filtrado con el elmento que conincide el search string. La UI todavia usa this.state.data para hacer  el renderizado pero this.state.data es un version reducida de ella misma
\\
Necesitas una referencia al dato antes de la busqueda, entonces no perdereas el dato para siempre. esto permite al usuario volver a la tabla completa o cambiar la busqueda para obtener diferentes coincidencias. Llamemos a esta referencia this.preSearchData. Ahora que hay datos guardados en dos lugares el metodo save() necesitara una actualizacion para ambos campos \\
Cuando el usuario hace click en el boton search la funcion toggleSearch() es invocada esta funcion prende el search y lo apaga.
\begin{itemize}
  \item Seteando el this.state.search en true o false
  \item Cuando es true, recuerda el dato actual
  \item Cuando es false revierte al dato 
\end{itemize}
\begin{lstlisting}[language=html]
toggleSearch(){
  if(this.state.search){
    this.setState({
      data: this.preSearchData,
      search: false,
    });
    this.preSearchData = null;
  } else {
    this.preSearchData = this.state.data;
    this.setState({
      search:true
    })
  }
  }
\end{lstlisting}
La ultima cosa para hacer es implementar el search, cual llama cada ves que algo es escrito en el search row onChange.
\begin{lstlisting}[language=html]
search(e){
  const needle = e.target.calue.toLowerCase();
  if (!needle){
    this.setState({data: this.preSearchData});
    return;
  }
  const idx = e.target.dataset.idx;
  const searchdata = this. preSearchData.filter((row) => {
    return row[idx].toString().toLowerCase().indexOf(needle) > -1;
  });
  this.setState({data: searchdata});
  }
\end{lstlisting}
obtienes la busqueda desde los cambios el evento
\section{Actualizar el metodo save()}
Se debe actulizar el metodo save para que soporte el this.preSearchData
\begin{lstlisting}[language=html]
save(e) {
  e.preventDefault();
  const input = e.target.firstChild;
  const data = clone(this.state.data).map((row) => {
    if (row[row.length - 1] === this.state.edit.row) {
      row[this.state.edit.column] = input.value;
    }
    return row;
  });
  this.logSetState({
    edit: null,
    data,
  });
  if (this.presearchData){
    this.preSearchData[this.state.edit.row][this.state.edit.column] = input.value;
  }
  }
\end{lstlisting}
\section{Como podrias mejorar el Search?}
Esto fue un simple trabajo de de ilustracion, puedes mejorar el feature para realizar un search de otro search?
\section{Instant Replay}
Como sabes tus componentes se preocupan del estado y React renderisa y renderisa donde sea apropiado. Esto signidifca que dado el mismo dato(state y properties) lucira exactamente igual\\
Imagina alguien encontrando un bug mientras usa la app, ellos pueden click un boton y reportar el bug sin necesidad de explicar que paso. este bug report puede solo enviarte atras una copia de this.state y this.props y tu seras capas de recrear el estado aplicacion exacta y ver el resultado visual.\\
un "undo" no seria otro feature, desde React renders your app la misma manera dado el mismo props y estado. es un hecho el "undo" implementacion es algo trivial: solo necsitas ir atras al estado previo\\
Tomemos una idea un poco exajerada por diversion. Gravemos cada cambio de estado en el Excel component y luego hagamos un replay. Es facinante ver todas las acciones reproduciendoce. La pregunta de cuando el cambio ocurre no es importante solo reproduscamos el estado de la app\\
Para implmentar este feature, necesitamos agregar un logSetState() cual primero logea el nuevo estado a un this.log array y despues llama setState() Donde sea en el codigo tu llamas setState() deberia ahora estar cambiada a call logSetState() primero, busca y remplasa todos las llamadas a setState() cual llama a la nueva funcion
\\
Todo lo que era
\\
this.setState()
\\
Ahora sera
\\
this.logSetState();
\\
Ahora movamonos al constructor. Necesitas bindear las dos nuevas funciones log SetState() y replay(), declara esto.log array, y asigna el estado inicial\\
\begin{lstlisting}[language=html]
constructor(props) {
  //....
  //log the initial state
  this.log = [clone(this.state)]
  this.replay = this.replay.bind(this)
  this.logSetState = this.logSetState.bind(this);
  }
\end{lstlisting}
El logSetState necesita hacer dos cosas: log el nuevo estado y despues pasarlo a setState() aqui un ejemplo de implementacion donde haces una cpia del estado
\begin{lstlisting}[language=html]
        logSetState(newState) {
          this.log.push(clone(newState));
          this.setState(newState);
        }
\end{lstlisting}
Ahora que todos los estados estan logueados reproduscamos de vuelta. Para lanzar el playback agregemos un evento simple listener que captura el keyboard action e invoca la funcion replay(). El lugar para event listener como esto es el metodo componenteDidMount() lifecycle:
\begin{lstlisting}[language=html]
        componentDidMount() {
          document.addEventListener("keydown", (e) => {
            if (e.altKey && e.shiftKey && e.keyCode === 82) {
              // Alt + SHIFT + R(eplay)
              this.replay();
            }
          });
        }
\end{lstlisting}
Finalmente considrando el metodo replay() es usado setInterval() y una vez por segundo esto lee el siguiente objeto desde el log y pasa a setState():
\begin{lstlisting}[language=html]
        replay() {
          if (this.log.length === 1) {
            console.warn("No state changes to replay yet");
            return;
          }
          let idx = -1;
          const interval = setInterval(() => {
            if (++idx === this.log.length - 1) {
              clearInterval(interval);
            }
            this.setState(this.log[idx]);
          }, 1000);
        }
\end{lstlisting}
Y con esto el nuevo feature esta completo puedes ver el ejemplo en 03.11.table.replay.html
\section{Limpiando manejador de Eventos}
El replay necesita un poco de limpieza. Cuando este componente es la unica cosa pasando en la pgina, limpar no es necseario, en una aplicacion real lo componentes son agregados y removidos del DOM frecuentenemente. Cuando removemos desde el DOM una buen ciudadano puede limpar despues de el mismo. En el ejemplo abajo hay dos items que necesitan ser limpiados: el keydown event listener y el replay interval callback\\
Si no limpias el keydown eventlistener fucntion, afectara a la memoria despues del componente se haya ido, Y por que esta usando this, toda la instancia Excell necesita ser retenida en meoria. Esto es memory leak, Muchas de estos y el usuario podria irse y la aplicacion podria crashear el browser tab. Como el intervalo bueno el callback continuara ejecutandoce despues del ocmponete se haya ido y cause otro memory leak. El callback tambien llamara a setState de un componente no existente\\
Puedes reproducir el error creando un hello world desde la consola
\begin{lstlisting}[language=html]
ReactDOM.render(
  React.createElement('h1', null, 'Hello world!'),
  document.getElementById('app'),
);
\end{lstlisting}
Tambien puedes logear un timestamp a la consola en el intervalo callback para ver que esto se mantiene ejecutando\\
\section{Cleaning Solution}
Tomando cuidado de este filtrado de memoria Necesitas mantener referencia al manejo e intervalo timeouts que quieres limpiar luego limpiarlo en componentWillUnmount() \\
Para el manejador evento  teniendo un class method opuesto a un inline function
\begin{lstlisting}[language=html]
        constructor(props) {
          // ...
          this.replayID = null;
          // ...
          this.keydownHandler = this.keydownHandler.bind(this);
        }
        // .......
        keydownHandler(e) {
          if (e.altKey && e.shiftKey && e.keyCode === 82) {
            // Alt + SHIFT + R(eplay)
            this.replay();
          }
        }
        componentDidMount() {
          document.addEventListener("keydown", this.keydownHandler);
          cleanInterval(this.replayID);
        }
        replay() {
          if (this.log.length === 1) {
            console.warn("No state changes to replay yet");
            return;
          }
          let idx = -1;
          this.replayID = setInterval(() => {
            if (++idx === this.log.length - 1) {
              clearInterval(this.replayID);
            }
            this.setState(this.log[idx]);
          }, 1000);
        }
\end{lstlisting}
\section{Puedes mejorar el Replay?}
Como implementar un undo/redo feature? cuando la persona usa el Alt + Z bas un paso atras en el stage log y Alt + Shift + Z vas al inicio
\section{Una Implementacion Alternativa}
Hay otra manera de implmentar replay/undo funcionality sin cambiar todo el setState()? Talvez usando una apropiado lifecycle method
\section{Download los datos de la tabla}
Despues de todos el sorting, edigint y searching el usuario esta finalmente feliz con el estado del dato en la tabla. Seria Genial si ellos podrian descargar los datos, el resultado\\
Afortunadamente, no hay nada facil en React Todo lo necesitas hacer es grabar el actual this.state.data y darlo de vuelta por ejemplo en formato JSON o (CSV) 
\\
La primera cosa a hacer es agregar nuevas opciones al toolbar(donde esta el Search button ) usamos alguna magia HTML para forsar <a> links para lanzar la descarga del archivo
\begin{lstlisting}[language=html]
        constructor(props){
          ......
          this.downloadJSON = this.download.bind(this, "json");
          this.downloadCSV = this.download.bind(this, "csv");
        }
        download(format, ev) {
          const data = clone(this.state.data).map((row) => {
            row.pop(); // drop the last column, the recordId
            return row;
          });
          const contents =
            format === "json"
              ? JSON.stringify(data, null, " ")
              : data.reduce((result, row) => {
                  return (
                    result +
                    row.reduce((rowcontent, cellcontent, idx) => {
                      const cell = cellcontent.replace(/"/g, '""');
                      const delimiter = idx < row.length - 1 ? "," : "";
                      return `\${rowcontent}"\${cellcontent}"\${delimiter}`;
                    }, "") +
                    "\n"
                  );
                }, "");
          const URL = window.URL || window.webkitURL;
          const blob = new Blob([contents], { type: "text/" + format });
          ev.target.href = URL.createObjectURL(blob);
          ev.target.download = "data." + format;
        }
        .....
              <div className="toolbar">
                <button onClick={this.toggleSearch}>
                  {this.state.search ? "Hidde search" : "Show search"}
                </button>
                <a href="data.json" onClick={this.downloadJSON}>
                  Export JSON
                </a>
                <a href="data.csv" onClick={this.downloadCSV}>
                  Export CSV
                </a>
              </div>
\end{lstlistign}
\section{Fetching Data}
El componente tiene acceso a los datos en el mimo archivo pero que si el dato esta en otro lugar, en un servidor, y necesitas ser obtenido? Hay varias soluciones para esto y veras mas opciones mas adleante, provemos una de las simples obteniendo datos en componentDidMount()\\
Digamos que el componente Excell es creado con un initialData property\\
El componente puede renderisar un estado intermedio para dejar saber al usuario que el dato esta en camino
\begin{lstlisting}[language=html]
       componentDidMount() {
          document.addEventListener("keydown", this.keydownHandler);
          fetch("https://www.phpied.com/files/reactbook/table-data.json")
            .then((response) => response.json())
            .then((initialData) => {
              const data = clone(initialData).map((row, idx) => {
                row.push(idx);
                return row;
              });
              this.setState({ data });
            });
          clearInterval(this.replayID);
        }
       .....
              {this.state.data.length === 0 ? (
                  <tbody>
                    <tr>
                      <td colSpan={this.props.headers.length}>
                        Loading data...
                    </td>
                  </tr>
                </tbody>
              ) : (
                <tbody onDoubleClick={this.showEditor}>
        .....
     ReactDOM.render(
        <Excell headers={headers} initialData={[]} />,
        document.getElementById("app"),
     );
\end{lstlisting}